\documentclass[../main]{subfiles}
\begin{document}
\section{Acuerdos áulicos}
\subsection*{Normas para la presentación de la carpeta de trabajos
prácticos, teoría y acuerdos áulicos.}
\begin{enumerate}
  \item Portada y formato general: La carpeta debe iniciar con la
    portada oficial de la escuela (o contener la misma información).
    Se exige el uso de letra clara y prolija. Se puede utilizar hoja
    común tamaño A4, pero es obligatorio que todas las hojas incluyan
    el nombre y apellido del alumno.
  \item Dinámica de clase: Las clases tienen modalidad de taller. Se
    trabajará con una guía de trabajos prácticos; primeramente se
    explicará el contenido teórico y luego los alumnos completarán la
    parte práctica durante la hora de clase.
  \item Uso de Google Classroom: Todos los alumnos deberán unirse al
    aula virtual. Los trabajos realizados se deben subir a la
    plataforma. Se evaluará tanto la correcta realización como el
    cumplimiento de la fecha de entrega pactada.
  \item Correcciones y reentregas: Si un trabajo práctico se entrega
    fuera de término o se encuentra incompleto, el alumno tendrá la
    obligación de corregirlo/completarlo y volver a enviarlo a través
    de Classroom.
  \item Revisión de carpeta física: Aunque las entregas regulares
    sean por Classroom, es obligatorio mantener la carpeta física
    ordenada y completa. A mitad de año (mes de junio,
    aproximadamente), el docente pedirá las carpetas para su revisión
    y corrección presencial.
  \item Proyecto grupal y evaluación continua: No habrá evaluaciones
    escritas tradicionales. En su lugar, los alumnos trabajarán en un
    proyecto grupal. En cada trabajo práctico habrá actividades
    específicas destinadas a nutrir este proyecto, las cuales deben
    ir actualizándose. Este trabajo se evaluará de forma continua,
    clase a clase.
  \item Cierre de la materia: La nota final integradora de la materia
    se definirá a fin de año mediante una exposición del proyecto
    frente al curso, que demostrará los conocimientos adquiridos a lo
    largo del ciclo lectivo.
  \item Pautas de convivencia y actitud: Se evaluará permanentemente
    la participación, el trabajo en grupo y en el pizarrón. El alumno
    deberá mantener un vocabulario adecuado, respetar al docente, a sus
    compañeros y cumplir con las normas de convivencia de la
    institución. En caso de inasistencia, se deberá presentar el
    certificado correspondiente al preceptor.
\end{enumerate}
\vspace{1cm}
\begin{center}
  Firma Alumno $\dots\dots\dots\dots\dots\dots\dots\dots\dots\dots$
\end{center}
\end{document}